%% kapa: exam=1 round=20 year=2009 period=1 form=A subject=economics subject_ko=경제원론 seq=11 number=51
%% source: Q-Net 감정평가사 기출문제 (cst00309) artlSeq=5024178 file=extracted/1008/5024178/2009년 제20회 감정평가사 1차시험 기출문제/1교시.hwp
%% conversion: kordoc 4.24.51 -> tools/gen1.py (auto); review=proofread-sonnet; flags=none
\begin{question}{51}
인플레이션율과 실업률에 관한 설명으로 옳은 것은?
\begin{choices}
  \item 필립스곡선은 인플레이션율과 실업률 간의 양(+)의 관계를 보여준다.
  \item 합리적 기대이론에 따르면 인플레이션율과 실업률은 장기적으로는 양(+)의 관계에 있지만 단기적으로는 음(-)의 관계에 있다.
  \item 인플레이션율과 실업률이 동시에 상승하는 현상은 필립스곡선의 이동으로도 설명될 수 있다.
  \item 스태그플레이션은 인플레이션율이 상승하면서 실업률이 감소되는 현상이다.
  \item 총공급 충격은 스태그플레이션을 초래할 수 없다.
\end{choices}
\end{question}
\answer{3}
